% 来源：解析据 papers/2016-2017-秋冬-期中-甲.tex 撰写（原卷见 sources/2016-2025-期中-甲-历年真题汇编.pdf 第 1 页）

\begin{solution}
第 $2$、$4$ 列只在第 $2$、$4$ 行有非零元，第 $1$、$3$、$5$ 列只在第 $1$、$3$、$5$ 行有非零元。把第 $2$、$4$ 行依次换到最前（$3$ 次行交换），再把第 $2$、$4$ 列依次换到最前（同样 $3$ 次列交换），两次的符号因子 $(-1)^{3}\cdot(-1)^{3}=1$ 相抵，行列式的值不变，而矩阵已化为分块对角形，故可用分块行列式：
\[
D=\begin{vmatrix}2&1\\5&3\end{vmatrix}\cdot
\begin{vmatrix}1&2&x\\1&4&x^{2}\\1&8&x^{3}\end{vmatrix}
=(6-5)\begin{vmatrix}1&2&x\\1&4&x^{2}\\1&8&x^{3}\end{vmatrix}.
\]
右端的 $3$ 阶行列式把第 $2$、$3$ 行分别减去第 $1$ 行，再按第 $1$ 列展开：
\[
\begin{vmatrix}1&2&x\\1&4&x^{2}\\1&8&x^{3}\end{vmatrix}
=\begin{vmatrix}1&2&x\\0&2&x^{2}-x\\0&6&x^{3}-x\end{vmatrix}
=\begin{vmatrix}2&x^{2}-x\\6&x^{3}-x\end{vmatrix}
=2(x^{3}-x)-6(x^{2}-x)=2x(x-1)(x-2).
\]
\ans{$D=2x(x-1)(x-2)=2x^{3}-6x^{2}+4x$。}
验算：$x=0$ 时第 $5$ 列全为零，$x=1$ 时第 $5$ 列与第 $1$ 列相同，$x=2$ 时第 $5$ 列与第 $3$ 列相同，故 $0,1,2$ 都是零点，与三次式相符。
\end{solution}

\begin{solution}
\textbf{先把方程组解清楚。}对增广矩阵作初等行变换（$R_2-R_1$，$R_3-4R_1$，$R_4-2R_1$，再 $R_3-3R_2$，$R_4+R_2$，最后 $R_4-\dfrac43R_3$）：
\[
\begin{pmatrix}
1&0&1&-1&-3&-2\\
1&2&-1&0&-1&1\\
4&6&-2&-4&3&7\\
2&-2&4&-7&4&1
\end{pmatrix}
\to
\begin{pmatrix}
1&0&1&-1&-3&-2\\
0&2&-2&1&2&3\\
0&0&0&-3&9&6\\
0&0&0&0&0&0
\end{pmatrix}.
\]
系数矩阵与增广矩阵的秩同为 $3$（主元在第 $1,2,4$ 列），故方程组相容，未知量有 $5$ 个，自由未知量有 $5-3=2$ 个。取 $x_3,x_5$ 为自由未知量，由化简后的三个非零行自下而上回代：
\begin{align*}
x_4&=3x_5-2,\\
x_2&=x_3+\frac{3-x_4-2x_5}{2}=x_3+\frac{5(1-x_5)}{2},\\
x_1&=-2-x_3+x_4+3x_5=6x_5-x_3-4 .
\end{align*}
直接令 $x_5=c_2$ 会出现分数，故取 $x_3=c_1$、$x_5=1+2c_2$（$x_5$ 仍取遍所有实数），得整系数的通解
\[
(x_1,x_2,x_3,x_4,x_5)^{T}=(2,0,0,1,1)^{T}+c_1(-1,1,1,0,0)^{T}+c_2(12,-5,0,6,2)^{T}.
\]
\ans{通解为 $\mathbf{x}=(2,0,0,1,1)^{T}+c_1(-1,1,1,0,0)^{T}+c_2(12,-5,0,6,2)^{T}$（$c_1,c_2$ 任意）。}
验算：取 $c_1=c_2=0$ 得 $(2,0,0,1,1)^{T}$，代入四个方程分别得 $-2,1,7,1$，与常数项一致；两个方向向量代入左端分别得零向量，故确为解。
\rem{题面写“齐次线性方程组”，但齐次方程组要求各方程的常数项全为零，而题面给出的方程常数项不全为零，故实际是非齐次方程组。这里按题面给出的方程（非齐次）求解。}
\end{solution}

\begin{solution}
\textbf{(1)} 在 $m\times n$ 矩阵 $A$ 中任取 $k$ 行与 $k$ 列，位于这些行、列交叉处的 $k^{2}$ 个元素按原次序排成的 $k$ 阶行列式，称为 $A$ 的一个 $k$ 阶子式。若 $A$ 中有一个 $r$ 阶子式不为零，而所有 $r+1$ 阶子式（若存在）全为零，则称 $A$ 的秩为 $r$，记作 $r(A)$；零矩阵的秩规定为 $0$。也就是说，$r(A)$ 就是 $A$ 中非零子式的最高阶数，它也等于 $A$ 的行阶梯形中非零行的行数。

\textbf{(2)} 对 $A$ 作初等行变换：先 $R_3\leftarrow R_3-2R_2$ 得 $(0,0,\lambda+2,1)$，再 $R_2\leftarrow R_2-R_1$ 得 $(0,\lambda+1,0,0)$，于是 $A$ 行等价于
\[
\begin{pmatrix}1&-1&1&1\\0&\lambda+1&0&0\\0&0&\lambda+2&1\end{pmatrix}.
\]
它至多有 $3$ 个非零行。当 $\lambda\ne-1$ 时第 $2$ 行非零：$\lambda\ne-2$ 时三行分别在第 $1,2,3$ 列取得主元，$\lambda=-2$ 时第 $3$ 行是 $(0,0,0,1)$、三行分别在第 $1,2,4$ 列取得主元。两种情形三行都线性无关，故 $r(A)=3$；当 $\lambda=-1$ 时第 $2$ 行全为零，剩下的 $(1,-1,1,1)$ 与 $(0,0,1,1)$ 线性无关，故 $r(A)=2$。
\ans{$r(A)=\begin{cases}3,&\lambda\ne-1,\\[2pt]2,&\lambda=-1.\end{cases}$}
\end{solution}

\begin{solution}
先判断可逆性。按第 $1$ 行展开得
\[
|A|=1\cdot(4-6)-2\cdot(8-2)+3\cdot(6-1)=-2-12+15=1\ne0 ,
\]
同理 $|B|=7\cdot5-9\cdot4=-1\ne0$，故 $A,B$ 都可逆。方程 $AXB=C$ 两边左乘 $A^{-1}$、右乘 $B^{-1}$，得
\[
X=A^{-1}CB^{-1}.
\]
用初等行变换求 $A^{-1}$，二阶的 $B^{-1}$ 由伴随矩阵公式 $B^{-1}=\dfrac{1}{|B|}B^{*}$ 直接写出：
\[
A^{-1}=\begin{pmatrix}-2&1&1\\-6&1&4\\5&-1&-3\end{pmatrix},
\qquad
B^{-1}=\frac{1}{-1}\begin{pmatrix}5&-9\\-4&7\end{pmatrix}
=\begin{pmatrix}-5&9\\4&-7\end{pmatrix}.
\]
先算 $A^{-1}C$（$3\times2$ 矩阵），再右乘 $B^{-1}$：
\[
A^{-1}C=\begin{pmatrix}1&-1\\3&0\\-2&1\end{pmatrix},
\qquad
X=\begin{pmatrix}1&-1\\3&0\\-2&1\end{pmatrix}
\begin{pmatrix}-5&9\\4&-7\end{pmatrix}
=\begin{pmatrix}-9&16\\-15&27\\14&-25\end{pmatrix}.
\]
验算：把 $X$ 代回，$AXB=\begin{pmatrix}1&2\\1&0\\2&3\end{pmatrix}=C$，成立。
\ans{$X=\begin{pmatrix}-9&16\\-15&27\\14&-25\end{pmatrix}$。}
\end{solution}

\begin{solution}
$A$ 是分块下三角矩阵，按第一行展开得 $|A|=2\cdot\bigl((-2)\cdot1-1\cdot(-1)\bigr)=-2\ne0$，且容易求出
\[
A^{-1}=\begin{pmatrix}\frac12&0&0\\0&-1&1\\0&-1&2\end{pmatrix},
\qquad
A^{*}=|A|A^{-1}=-2A^{-1}=\begin{pmatrix}-1&0&0\\0&2&-2\\0&2&-4\end{pmatrix}.
\]
关键一步是用 $|A|=-2$ 把 $A^{*}$ 换成 $-2A^{-1}$，于是原方程化为
\[
(-2A^{-1}-2E)BA=-12E
\quad\Longrightarrow\quad
(A^{-1}+E)BA=6E .
\]
两边右乘 $A^{-1}$，得 $(A^{-1}+E)B=6A^{-1}$，即 $B=6(A^{-1}+E)^{-1}A^{-1}$。又 $(A^{-1}+E)A=A^{-1}A+A=E+A$，由 $(AQ)^{-1}=Q^{-1}A^{-1}$（取 $Q=A^{-1}+E$）得
\[
(A^{-1}+E)^{-1}A^{-1}=\bigl(A(A^{-1}+E)\bigr)^{-1}=(A+E)^{-1},
\]
故 $B=6(A+E)^{-1}$。而
\[
A+E=\begin{pmatrix}3&0&0\\0&-1&1\\0&-1&2\end{pmatrix},
\qquad
(A+E)^{-1}=\begin{pmatrix}\frac13&0&0\\0&-2&1\\0&-1&1\end{pmatrix},
\]
于是 $B=6(A+E)^{-1}$。
验算：代入得
\[
BA=\begin{pmatrix}4&0&0\\0&18&-6\\0&6&0\end{pmatrix},
\qquad
A^{*}BA=2BA-12E=\begin{pmatrix}-4&0&0\\0&24&-12\\0&12&-12\end{pmatrix},
\]
成立。
\ans{$B=6(A+E)^{-1}=\begin{pmatrix}2&0&0\\0&-12&6\\0&-6&6\end{pmatrix}$。}
\rem{伴随矩阵与转置容易混淆，这里不妨验一下：若把 $A^{*}$ 误当转置，方程化为 $(A^{T}-2E)BA=-12E$，而 $A$ 的首行与首列都是 $(2,0,0)$，故 $A^{T}-2E$ 的首行全为零，左端乘积的首行必为零行，右端 $-12E$ 的首行却是 $(-12,0,0)$，矛盾，方程无解。可见此处的 $A^{*}$ 只能是伴随矩阵。}
\end{solution}

\begin{solution}
记 $M=\begin{pmatrix}O&A\\B&C\end{pmatrix}$，其中 $O$ 是 $r\times s$ 零块、$C$ 是 $s\times r$ 块，故 $M$ 是 $r+s$ 阶方阵。按“先猜后验”的办法，把待求的逆矩阵按块行 $(s,r)$、块列 $(r,s)$ 分块（即左上 $s\times r$、右下 $r\times s$），取
\[
X=\begin{pmatrix}-B^{-1}CA^{-1}&B^{-1}\\A^{-1}&O\end{pmatrix},
\]
用分块乘法逐块检验：
\[
MX=
\begin{pmatrix}
O(-B^{-1}CA^{-1})+AA^{-1} & OB^{-1}+AO\\[2pt]
B(-B^{-1}CA^{-1})+CA^{-1} & BB^{-1}+CO
\end{pmatrix}
=
\begin{pmatrix}E_r&O\\O&E_s\end{pmatrix}
=E_{r+s},
\]
其中左下块 $B(-B^{-1}CA^{-1})+CA^{-1}=-CA^{-1}+CA^{-1}=O$，右上块 $OB^{-1}+AO=O$；各行用到了 $A,B$ 可逆给出的 $AA^{-1}=E_r$、$BB^{-1}=E_s$。于是 $M$ 可逆，且 $M^{-1}=X$。

若只想先判断可逆，也可算行列式：把两块行组互换需 $r\cdot s$ 次行交换，互换后成为分块三角阵 $\begin{pmatrix}B&C\\O&A\end{pmatrix}$，故 $|M|=(-1)^{rs}|A||B|\ne0$。
\ans{$\begin{pmatrix}O&A\\B&C\end{pmatrix}^{-1}=\begin{pmatrix}-B^{-1}CA^{-1}&B^{-1}\\A^{-1}&O\end{pmatrix}$。}
\end{solution}

\begin{solution}
\textbf{(1)} 设 $A$ 是 $m\times n$ 矩阵，其行向量为 $\alpha_1^{T},\alpha_2^{T},\dots,\alpha_m^{T}$。先记一条乘法事实：对任意 $m$ 阶矩阵 $P=(p_{ij})$，由 $(PA)_{ik}=\sum_{j=1}^{m}p_{ij}a_{jk}$ 得
\[
PA\ \text{的第 } i\ \text{行}=\sum_{j=1}^{m}p_{ij}\alpha_j^{T},
\]
即 $P$ 的第 $i$ 行元素正是把 $A$ 的各行作线性组合时的系数。

对 $m$ 阶单位矩阵 $E$ 作一次初等行变换，得到的矩阵称为相应的初等矩阵，共三类，逐类验证即可：
\begin{enumerate}[label=(\roman*),leftmargin=2.4em,itemsep=0.2em,topsep=0.25em]
  \item 互换第 $i,j$ 行得 $E(i,j)$，它的第 $i$ 行是 $e_j^{T}$、第 $j$ 行是 $e_i^{T}$，其余第 $k$ 行仍为 $e_k^{T}$。于是 $E(i,j)A$ 的第 $i$ 行是 $\alpha_j^{T}$、第 $j$ 行是 $\alpha_i^{T}$，其余行不变，恰为互换 $A$ 的第 $i,j$ 行。
  \item 第 $i$ 行乘非零数 $k$ 得 $E(i(k))$，它的第 $i$ 行是 $ke_i^{T}$，于是 $E(i(k))A$ 的第 $i$ 行是 $k\alpha_i^{T}$，其余行不变，恰为把 $A$ 的第 $i$ 行乘 $k$。
  \item 第 $j$ 行的 $k$ 倍加到第 $i$ 行得 $E(i,j(k))$，它的第 $i$ 行是 $e_i^{T}+ke_j^{T}$，于是 $E(i,j(k))A$ 的第 $i$ 行是 $\alpha_i^{T}+k\alpha_j^{T}$，其余行不变，恰为把 $A$ 的第 $j$ 行的 $k$ 倍加到第 $i$ 行。
\end{enumerate}
三类初等行变换都如此，故对 $A$ 施行一次初等行变换，等价于用相应的初等矩阵左乘 $A$。顺带由上还可读出：初等矩阵都可逆，且逆仍为同类初等矩阵（$E(i,j)^{-1}=E(i,j)$，$E(i(k))^{-1}=E(i(1/k))$，$E(i,j(k))^{-1}=E(i,j(-k))$）。

\textbf{(2)} 必要性：设 $A$ 可逆，则 $AA^{-1}=E$。两边取行列式并用 $|XY|=|X||Y|$，得 $|A|\cdot|A^{-1}|=1$，故 $|A|\ne0$（且 $|A^{-1}|=\dfrac{1}{|A|}$）。

充分性：设 $|A|\ne0$。用初等行变换把 $A$ 化为行阶梯形 $U$。由 (1)，$U=PA$，其中 $P$ 是有限个初等矩阵之积；三类初等矩阵的行列式分别为 $-1,k\ (k\ne0),1$，都非零，故 $|P|\ne0$，从而
\[
|U|=|P|\cdot|A|\ne0 .
\]
特别地 $U$ 没有零行。$n$ 阶矩阵的行阶梯形没有零行，说明它有 $n$ 个主元，即 $U$ 是主元全不为零的上三角矩阵。再对 $U$ 作倍加变换（把第 $i$ 行减去上方各行的适当倍数）消去每个主元上方的元素，并用倍乘变换把各主元化为 $1$，就把 $U$ 化成了 $E$。即 $A$ 行等价于 $E$，故存在初等矩阵之积 $Q$ 使 $QA=E$。

由 $QA=E$ 两边左乘 $Q^{-1}$ 得 $A=Q^{-1}$；而由 (1) 后的说明 $Q$ 可逆，故 $A$ 可逆（且 $A^{-1}=Q$，即 $A^{-1}$ 也是初等矩阵之积）。必要性、充分性都成立。
\qed
\end{solution}
